\chapter{Related work}\label{chap:related-work}
CTaps is not the only implementation of TAPS, as several
related open and closed source projects exist. This chapter
will describe the current state of the existing TAPS implementations,
forming a basis for why a new open source implementation is still
valuable. It will also compare the features currently present in CTaps with
the most complete previous implementations of TAPS, highlighting strengths
and weaknesses of CTaps.

\section{NEAT}
One of the first open source implementations of TAPS is 
NEAT~\cite{khademiNEATPlatformProtocolIndependent2017, neatCode},
whose development was funded by the EU horizon 2020 project
from 2015 to 2018~\cite{neatFunding}.
Developed alongside TAPS, the development of TAPS and NEAT naturally
affected each other. However, while NEAT halted development in 2018,
the TAPS API has continued to evolve, and the final TAPS RFC~was published
in 2025~\cite{rfc9622}.

NEAT is an asynchronous library and
uses an event loop, allowing for users to add their own callbacks.

The core architecture of TAPS and NEAT are similar, with both being
asynchronous, event-driven APIs, where the underlying connections
can be configured through properties.

Within the documented properties, NEAT has a different approach
to central concepts than TAPS. For example with the Transport Property,
the application can specify reliable, but also directly
specify one or more potential transport protocols to use. This
is in contrast to TAPS, where the application only interfaces
with the protocol through selection properties.

\section{NEATPy}
NEATPy~\cite{gundersenNEATPyStandardsConformantProtocolIndependent2020}
is a Python library developed as a master's thesis at UiO in 2020.
It was developed as an attempt to create a
Python layer to bridge the growing gap between NEAT and TAPS. Among
the goals was to move away from specifying the protocol at design time
instead using Selection Properties as specified by TAPS.

Because NEATPy is a Python shim on top of NEAT, certain parts of it are
limited by NEAT itself. For example,
using NEATPy it is not possible to specify the \texttt{Interface Instance or Type}
or the \texttt{Use Temporary Local Address} Selection Properties 
\cite{NEATPyDocumentationNEATPy}. In addition, NEATPy does not
support QUIC, as this is not supported by NEAT. 

Another
limitation of NEATPy is that it was developed before the finalization of
the TAPS standard. Developed in 2020, NEATPy is based mostly on
version 4 and 5 of the TAPS interface draft~\cite[p. 109]{gundersenNEATPyStandardsConformantProtocolIndependent2020}.
The TAPS API has since evolved, and was 
published as a proposed standard in January 2025
\cite{rfc9622}.

Since the finalization of NEATPy which implements
version 5 of the TAPS draft~\cite{ietf-taps-interface-05} several
new features have been added to TAPS.
For example, the support for giving more than one local and remote
endpoint when creating a Preconnection. This is done when the application
wants to indicate that several endpoints are eligible for the connection
or provide the same service.

In addition, there are several new Selection Properties that have been added
to TAPS. These are:
\begin{itemize}
    \item Keep-Alive Packets
    \item Use Temporary Local Address
    \item Multipath Transport
    \item Advertisement of Alternative Addresses
    \item Initiating Side Is Not the First to Write
\end{itemize}

Some selection properties have also been removed from the TAPS API:
\begin{itemize}
    \item Parallel Use of Multiple Paths
    \item Notification of excessive retransmissions
\end{itemize}

In addition to the outward differences to the API, there are also several
changes to the required internal workings of TAPS since the development of NEATPy.
For example, in the final RFC~TAPS has loosened the required
order of operations when working with selection properties.
Previously, there was a strict precedence in the order of
which preference values were to be applied.
This has since been
relaxed.% and the prioritization has been moved from being
% related to the preference level, to being ordered by which
% selection properties will affect the path selection.

\section{Network.framework}
Apple's Network.framework~\cite{networkFramework} follows some of the same concepts
laid out by TAPS, although it is not a complete implementation. It
features several of the same abstraction objects
such as Connections and Listeners and presents itself as an alternative 
to using sockets directly~\cite{wwdc2018}. There are however several key differences between
Network.framework and a fully fledged TAPS system, with the most important
being the lack of Selection Properties and racing of protocol candidates.
\todo{Cite that repo where apple says its a taps impl.}

With Network.framework, the underlying transport protocol is specified in
the constructor of the Connection object, and while it does support
Happy Eyeballs-style racing of IPv4/IPv6 addresses, it does not support
racing across multiple remote endpoints or across different transport protocols, as TAPS describes.

\section{PyTAPS}
PyTAPS~\cite{frankePythonasynciotaps2025} is a Python library aiming
to be fully compliant with TAPS. Unlike NEATPy, PyTAPS is not built
on top of NEAT, and therefore does not suffer from the same limitations
that NEATPy does. However, PyTAPS is based on the fourth version of the TAPS
interface draft~\cite{frankePythonasynciotaps2025} and has not had active
development since 2020. As such, it naturally suffers from the same
issues as NEATPy when it comes to the changes in the TAPS API.

\section{Dynamic TAPS}
dynamic-taps is a C implementation of TAPS, developed by the
American technology company F5~\cite{dukeDynamictaps2025}. It
is in the very early stages of development and has not had
active development since 2021, several years before the
finalzation of the TAPS RFCs.

\section{Comparison of TAPS implementations}\label{sec:tapsfeatures}
\todo{Discuss and finish table}

\begin{table}[htbp]
\centering
\begin{threeparttable}
\caption{Comparison of features of various TAPS implementations}
\label{tab:comparison}
\small
\begin{tabular}{lccccc}
\toprule
\rowcolor{gray!25}
\textbf{Feature}                                     & \textbf{NEATPy\tnote{1}} & \textbf{PyTAPS} & \textbf{Network.framework}     & \textbf{CTaps}\\
\midrule
\multicolumn{1}{l||}{Post TAPS release\tnote{2}}    & \xmark                   & \xmark          & \cmark                     & \cmark         \\
\multicolumn{1}{l||}{Language}                       & Python                   & Python          & Swift / Obj.-C              & C         \\
\multicolumn{1}{l||}{Candidate Racing}               & \cmark                   & \cmark          & \xmark\tnote{3}            & \cmark         \\
\multicolumn{1}{l||}{TLS support}                    & \cmark                   & \xmark          & \cmark                     & \xmark         \\
\multicolumn{1}{l||}{QUIC support}                   & \xmark                   & \xmark          & \cmark                     & \cmark         \\
\multicolumn{1}{l||}{Multipath\tnote{4}}             & \xmark                   & \xmark          & \cmark\tnote{5}            & \cmark         \\
\multicolumn{1}{l||}{Early data}                     & \xmark                   & \xmark          & \cmark\tnote{6}            & \cmark         \\
\multicolumn{1}{l||}{Multistreaming}                 & \xmark                   & \xmark          & \cmark                     & \cmark         \\
\multicolumn{1}{l||}{Rendezvous}                     & \xmark                   & \xmark          & \cmark                     & \xmark         \\
\multicolumn{1}{l||}{OS support\tnote{7}}            & B, L, M                  & L, M\tnote{8}   & M                          & L         \\
\multicolumn{1}{l||}{Open source}                    & \cmark                   & \cmark          & \xmark                     & \cmark         \\
\bottomrule
\end{tabular}
\begin{tablenotes}
    \footnotesize
    \item[1] NEAT is excluded as its features would be a subset of NEATPy's.
    \item[2] Whether the project has been under active development after the full release of the TAPS RFCs in 2025.
    \item[3] Supports Happy Eyeballs-style IP family racing~\cite{networkframework_happy_eyeballs_wwdc}, but not full candidate racing as described by TAPS.
    \item[4] Taken to mean that a single connection can use more than one path in a single session,
        either concurrently or through connection migration.
    \item[5] Supports multipath TCP~\cite{apple_multipathservicetype}, QUIC support is unclear.
    \item[6] Supported for TCP~\cite{apple_allowfastopen}.
        Not supported in QUIC according to an Apple Engineer's post on the Apple Developer Forums in 2023~\cite{apple_dts_0rtt_2023}.
        Lacks explicit support in the NWProtocolQUIC.Options class~\cite{apple_nwprotocolquic}.
    \item[7] B = BSD, L = Linux, M = macOS, W = Windows.
    \item[8] Based on PyTAPS' README
\end{tablenotes}
\end{threeparttable}
\end{table}
